\section{THỰC TRẠNG HOẠT ĐỘNG QUẢN LÝ VÀ MÔ HÌNH QUẢN LÝ TẠI FPT}

\subsection{Bộ máy điều hành và các nhân vật lãnh đạo cấp cao}
Nếu quản trị là việc định ra luật chơi và giám sát cuộc chơi, thì \textbf{Quản lý (Management)} là quá trình trực tiếp điều khiển các nguồn lực (con người, tài chính, công nghệ, thông tin) để hoàn thành các mục tiêu đề ra.

Ban Điều hành là bộ máy thực thi của FPT, chịu trách nhiệm tổ chức vận hành hoạt động sản xuất kinh doanh thường nhật dưới sự ủy quyền và giám sát của HĐQT.

\subsubsection{Lãnh đạo cấp tập đoàn}
Đội ngũ Ban Điều hành cấp Tập đoàn FPT gồm những nhà quản lý chuyên nghiệp, gắn bó lâu năm và am hiểu sâu sắc thị trường CNTT (Bảng \ref{tab:ban_dieu_hanh}):

\begin{table}[htbp]
\centering
\caption{Nhân sự Ban Điều hành cấp Tập đoàn FPT}
\label{tab:ban_dieu_hanh}
\begin{tabularx}{\textwidth}{p{4.2cm} p{4.0cm} X}
\toprule
\textbf{Họ và tên} & \textbf{Chức vụ} & \textbf{Trách nhiệm và phân công quản lý} \\
\midrule
\textbf{Ông Nguyễn Văn Khoa} & Tổng Giám đốc (CEO) & Bổ nhiệm từ năm 2019; điều hành toàn diện hoạt động kinh doanh, thực thi chiến lược chuyển đổi số và phát triển thị trường toàn cầu. \\
\textbf{Ông Nguyễn Thế Phương} & Phó Tổng Giám đốc & Đảm nhiệm vị trí từ năm 2000; phụ trách công tác tài chính, kế hoạch kế toán và các khối nghiệp vụ tập trung. \\
\textbf{Ông Phạm Minh Tuấn} & Phó Tổng Giám đốc & Bổ nhiệm năm 2024; đồng thời là CEO FPT Software, điều phối chiến lược mở rộng thị trường phần mềm quốc tế. \\
\textbf{Ông Hoàng Hữu Chiến} & Kế toán trưởng & Chịu trách nhiệm về hệ thống kế toán, lập báo cáo tài chính hợp nhất và kiểm soát dòng tiền tập đoàn. \\
\bottomrule
\end{tabularx}
\end{table}

\subsubsection{Đội ngũ lãnh đạo cấp cao tại các đơn vị thành viên}
FPT sở hữu đội ngũ quản lý cấp cao tinh nhuệ tại các mảng kinh doanh trọng yếu:
\begin{itemize}
    \item \textbf{Ông Phạm Minh Tuấn:} Phó Tổng Giám đốc FPT kiêm Tổng Giám đốc FPT Software -- chỉ huy mũi nhọn xuất khẩu phần mềm với doanh thu vượt mốc 1 tỷ USD.
    \item \textbf{Ông Nguyễn Khải Hoàn:} Phó Tổng Giám đốc kiêm Giám đốc Vận hành (COO) FPT Software, Tổng Giám đốc FPT Americas -- phụ trách thị trường trọng điểm châu Mỹ.
    \item \textbf{Ông Đỗ Văn Khắc:} Lãnh đạo cấp cao FPT Software, Tổng Giám đốc FPT Asia Pacific -- chỉ huy thị trường châu Á -- Thái Bình Dương.
    \item \textbf{Ông Ngô Mạnh Cường:} Nguyên CEO FPT Online, được bổ nhiệm giữ chức Phó Tổng Giám đốc FPT Smart Cloud từ ngày 21/9/2026.
    \item \textbf{Ông Trần Hải Linh:} Nguyên CEO Sendo, được bổ nhiệm giữ chức Tổng Giám đốc FPT Online từ ngày 21/9/2026 nhằm thúc đẩy đổi mới truyền thông số.
\end{itemize}

\subsection{Hệ thống các cấp quản lý trong tổ chức}
Để quản trị và điều phối đội ngũ hơn 70.000 cán bộ nhân viên toàn cầu, FPT phân chia thành 7 cấp bậc quản lý mạch lạc từ trên xuống dưới:
\begin{enumerate}
    \item \textbf{Cấp 1 -- Ban Điều hành Tập đoàn (C-Level Tập đoàn):} CEO, các Phó TGĐ và Kế toán trưởng -- thiết lập chỉ tiêu kinh doanh và phân bổ nguồn lực chung.
    \item \textbf{Cấp 2 -- CEO / President các công ty thành viên:} Người đứng đầu FPT Software, FPT IS, FPT Telecom, FPT Smart Cloud, FPT Education... chịu trách nhiệm toàn diện về kết quả P\&L (lãi/lỗ) của đơn vị.
    \item \textbf{Cấp 3 -- Phó Tổng Giám đốc công ty thành viên:} Phụ trách từng phân mảng chức năng chuyên sâu như Vận hành (COO), Công nghệ (CTO), Kinh doanh (CCO), Nhân sự (CHRO).
    \item \textbf{Cấp 4 -- Giám đốc Khối chức năng và Giám đốc Thị trường (Country Manager):} Quản lý các thị trường quốc tế trọng điểm (FPT Japan, FPT USA, FPT Europe...) hoặc các khối quản trị tập trung (Tài chính, Pháp chế, Bảo mật, Đào tạo).
    \item \textbf{Cấp 5 -- Giám đốc Sản xuất / Quản lý Chuyển giao dự án (Delivery Manager / PMO):} Điều phối các trung tâm sản xuất phần mềm (Software Development Center), chịu trách nhiệm về ngân sách, tiến độ bàn giao và chất lượng dự án công nghệ.
    \item \textbf{Cấp 6 -- Trưởng nhóm dự án / Kỹ thuật (Project Manager / Tech Lead / Scrum Master):} Trực tiếp chỉ đạo các nhóm kỹ sư từ 10 đến 50 người, phân chia công việc hàng ngày, giám sát quy trình Agile/Scrum.
    \item \textbf{Cấp 7 -- Nhân viên tác nghiệp trực tiếp:} Lập trình viên, kỹ sư kiểm thử (Tester), chuyên viên phân tích nghiệp vụ (BA), nhân viên kinh doanh, chuyên viên hỗ trợ khách hàng, vận hành hạ tầng...
\end{enumerate}

\subsection{Xác định mô hình quản lý của FPT}
Về mặt lý luận khoa học quản lý, cơ cấu quản lý của FPT không đơn thuần là một mô hình truyền thống mà là sự tổng hòa của: \textbf{Mô hình quản lý Công ty Mẹ -- Con (Holding Structure) kết hợp với Cơ cấu Ma trận đa chiều (Multidimensional Matrix Structure)}.

\begin{figure}[htbp]
\centering
\begin{verbatim}
                 BAN TỔNG GIÁM ĐỐC TẬP ĐOÀN
                             |
         +-------------------+-------------------+
         |                                       |
  CÁC KHỐI CHỨC NĂNG                 CÁC ĐƠN VỊ KINH DOANH
  (HR, Tài chính, Pháp chế, AI)      CHIẾN LƯỢC (SBU)
         |                                       |
         +-------------------X-------------------+
                             |
             CƠ CẤU MA TRẬN TẠI CÁC THỊ TRƯỜNG:
   - FPT Software Japan / Americas / Europe / APAC
   - FPT IS Tài chính công / Doanh nghiệp / Y tế
   - Dự án: Delivery Manager <---> Resource Manager
\end{verbatim}
\caption{Mô hình quản lý ma trận kết hợp SBU tại Tập đoàn FPT}
\label{fig:mo_hinh_quan_ly}
\end{figure}

\subsubsection{Mô hình Công ty Mẹ -- Con (Holding Structure)}
FPT Corporation đóng vai trò là công ty mẹ nắm giữ quyền kiểm soát vốn tại các công ty con và công ty liên kết. Mô hình này mang lại các đặc trưng:
\begin{itemize}
    \item \textbf{Tập trung hóa quyền lực vĩ mô:} Công ty mẹ nắm quyền hoạch định chiến lược dài hạn, phê chuẩn ngân sách đầu tư lớn, quản trị rủi ro tập trung, gìn giữ văn hóa cốt lõi và thương hiệu \textit{``FPT''}.
    \item \textbf{Phi tập trung hóa quyền điều hành tác nghiệp:} Các công ty con (FPT Software, FPT IS...) có tư cách pháp nhân riêng, được trao quyền tự chủ tối đa trong việc tìm kiếm khách hàng, tổ chức sản xuất và hạch toán kinh doanh. Điều này giúp các đơn vị thích ứng mau lẹ với đặc thù từng thị trường.
\end{itemize}

\subsubsection{Cơ cấu Ma trận đa chiều (Matrix Structure)}
Trong từng đơn vị thành viên (tiêu biểu là FPT Software và FPT IS), mô hình ma trận được áp dụng triệt để:
\begin{itemize}
    \item \textbf{Giao thoa giữa Khối sản phẩm (Business Unit) và Khối địa bàn (Geographical Unit):} Một kỹ sư hoặc nhóm dự án có thể phục vụ khách hàng thuộc chi nhánh FPT Japan (khối địa bàn), đồng thời trực thuộc Trung tâm năng lực công nghệ AI/Cloud (khối chuyên môn kỹ thuật).
    \item \textbf{Cơ chế hai tuyến báo cáo (Dual Reporting):} Trong các dự án phần mềm, nhân viên vừa báo cáo trực tiếp cho Giám đốc Dự án (Project Manager) về tiến độ, chất lượng sản phẩm công việc; vừa báo cáo cho Trưởng bộ phận chức năng (Resource Manager / Group Leader) về lộ trình thăng tiến, đánh giá năng lực và quyền lợi đãi ngộ.
\end{itemize}

\subsection{Cơ chế vận hành và công cụ phối hợp điều hành}
Để điều hành một bộ máy khổng lồ hoạt động trơn tru, FPT áp dụng các công cụ quản lý hiện đại:
\begin{itemize}
    \item \textbf{Hệ thống quản trị mục tiêu OKR (Objectives and Key Results):} Định kỳ hàng quý, các mục tiêu lớn đầy thách thức của Tập đoàn được phân rã thành các kết quả then chốt đo lường được cho từng công ty thành viên, phòng ban và đến từng cá nhân nhân viên.
    \item \textbf{Bộ chỉ số hiệu suất BSC/KPI (Balanced Scorecard):} Cân bằng giữa 4 khía cạnh: Tài chính, Khách hàng, Quy trình nội bộ và Học hỏi -- Phát triển.
    \item \textbf{Quy trình chuẩn hóa quốc tế:} Áp dụng tiêu chuẩn quản lý chất lượng ISO 9001 và mô hình trưởng thành năng lực CMMi Level 5 trong sản xuất phần mềm, đảm bảo tỷ lệ lỗi ở mức thấp nhất khi bàn giao sản phẩm cho khách hàng quốc tế.
\end{itemize}
